\section{Discussion}
\label{sec:discussion}
% \devroop{go over this and rewrite as needed} -- SOLVED: Rewrote the Discussion as a concise synthesis of the validated Results claims and removed rendered audit text.

Across corpus and external testbeds, pursuit-based hybrids consistently define the strongest efficiency--stability frontier, while adversarial-first baselines and poorly matched allocators expose lower robustness floors under structured and adaptive disruption. Our results show that routing robustness under compound uncertainty is jointly determined by \emph{policy context}, \emph{allocator choice}, and \emph{replay-capacity semantics}.
% The appendix provides further details regarding these summarized findings. \Cref{tab:testbed_comparison_full,tab:model_family_comparison_full} report the per-configuration cross-testbed and model-family results supporting the robustness hierarchy, while \Cref{fig:appendix_context_vs_exp3_gap,fig:appendix_context_capacity_effects,fig:appendix_threat_rules,fig:appendix_g8_advanced_synthesis} give additional diagnostics for Oracle-gap behavior, context-capacity behavior, deployment rules, replay sensitivity, and grouped claim interactions.
The appendix \Cref{tab:model_family_comparison_full} shows a 1--3 pp separation between \texttt{iCPursuitNeuralUCB} and the nearest internal neural alternatives. \Cref{tab:testbed_comparison_full} shows that its average-efficiency lead persists across the four evaluated external testbeds, although \texttt{EXPNeuralUCB} wins more individual configurations on \paperEight; average ranking and configuration-level wins can therefore diverge as topology and settings change.
\Cref{fig:appendix_context_deployment_support,fig:appendix_g8_advanced_synthesis} provide the corresponding visual audit: the first figure connects Oracle-gap, qubit-budget, and policy-ranking diagnostics, while the second groups replay drop/recovery, allocator, context-capacity, capacity-trajectory, and regret evidence behind the deployment interpretation.

% The appendix validates neural hybrid performance in detail. Neural hybrids achieve 
% 87--96\% efficiency with CV≤6.5\% \Cref{tab:model_family_comparison_full} and 
% outperform non-contextual baselines by 18--24 pp \Cref{tab:testbed_comparison_full}, 
% confirming context-awareness provides substantial advantage. Under attacks, they sustain 
% 5--10 pp Oracle gaps versus 20--40 pp for adversarial-first methods 
% \Cref{fig:appendix_context_vs_exp3_gap}, defining the stability frontier. The capacity 
% paradox—threat-conditioned replay sensitivity in \Cref{fig:appendix_g8_advanced_synthesis}—shows
% predictability, not bandwidth, limits robustness. \Cref{fig:appendix_threat_rules} maps 
% threat-responsive allocator rules, while \Cref{fig:appendix_g8_advanced_synthesis} 
% synthesizes how algorithm, allocator, and capacity jointly determine worst-case performance.

% ----------------------------------------------------------------------------
\smallTitle{Context-Aware Bandit Policies vs the Robustness Frontier}

Across the paper baselines (\RQOneText--\RQTwoText) and the Hybrid suite (\RQThreeText), context-aware policies outperform non-contextual and purely adversarial baselines in the evaluated threat settings. Pursuit-based contextual hybrids (\texttt{CPursuitNeuralUCB}, \texttt{iCPursuitNeuralUCB}) maintain strong efficiency across the five threat regimes, while EXP3-family baselines show weaker robustness floors under Markov, Adaptive, and OnlineAdaptive disruption. Within this simulator evidence, topology and channel-state context are associated with stronger routing robustness.

% ----------------------------------------------------------------------------
\smallTitle{The Capacity Paradox is Conditional}

Increasing replay capacity can improve stability in structured regimes while amplifying vulnerability under reactive threats. The observed interaction is consistent with allocator-driven exploration predictability, but the current experiments do not isolate that mechanism. Capacity sensitivity is most severe under allocator stress, particularly with \texttt{Random}; in the Hybrid suite, efficiency swings exceed 20--38 percentage points across replay-capacity settings under Adaptive and Markov threats. Replay-capacity effects must therefore be interpreted jointly with allocator policy and threat regime.

% ----------------------------------------------------------------------------
\smallTitle{Allocator Choice is the Practical Deployment Lever}

Allocator choice produces the largest observed configuration shifts in the Hybrid suite. Aggregated across scenarios and replay-capacity settings, \texttt{ThompsonSampling} and \texttt{Fixed}/\texttt{Default} form the top performance tier, while \texttt{Random} is associated with substantially lower efficiency floors. Within the evaluated conditions, allocator choice is therefore a first-class robustness control: even strong pursuit-neural architectures enter the capacity paradox under allocator mismatch, while the higher-performing allocators reduce replay sensitivity.

% ----------------------------------------------------------------------------
\subsection{Implications for Adversarial Quantum Network Design}

Within the evaluated settings, the results support three configuration principles: include contextual routing features, treat replay capacity as a coupled control variable, and adapt settings to the observed threat regime. Topology, link-state, and traffic context produce larger robustness gains than reward-only exploration in the tested comparisons; replay anchoring, scale, and allocator jointly shape tail risk under structured and adaptive disruption. \texttt{ThompsonSampling} and the static \texttt{Default} allocator are therefore candidate conservative baselines for future higher-fidelity evaluation, not hardware-validated deployment prescriptions.

% First, context-aware routing should be treated as baseline infrastructure. Feature enrichment through topology, link-state, and traffic information yields larger robustness gains than refining reward-only exploration rules.

% Second, replay capacity must be treated as a coupled control variable. The relevant deployment question is not simply \textquotedblleft more or less replay,\textquotedblright but which anchoring rule (\small\texttt{$T_b$} vs.\ \small\texttt{$T$}), scale ($s$), and allocator jointly minimize tail risk under the active threat regime.

% Finally, the results motivate lightweight threat-aware meta-policies:
% \begin{itemize}[nosep]
% \item Use \small\texttt{ThompsonSampling} or \small\texttt{Default} allocators as safe deployment defaults.
% \item Treat replay-capacity settings as \emph{scenario-coupled}.
% \item Dynamically adapt allocator and replay configurations using efficiency and variability signals observed during operation.
% \end{itemize}
